\subsection{What the class-level errors reveal}
The standalone image classifier's mean observed-test F1 ranges from 85.3\% for pedestrians to 0.0\% for taxis. The coordinate-based family reaches 27.8\% on taxis, still a weak result. Across the three diffusion seeds the mean number of distinct images in an eight-candidate test pool is 8.00, 7.98 and 8.00. Thus the scorer typically sees distinct inputs, even though distinct images need not give useful distinctions among the requested classes.

A second view is to inspect downstream recall. For the TCN, adding GAF-selected examples increases bus recall from 84.6\% to 92.3\%, while motorcycle recall decreases from 51.6\% to 43.0\% and taxi recall from 62.5\% to 45.8\%. For the GRU, bus recall rises from 41.0\% to 59.0\%, while motorcycle recall decreases from 37.6\% to 30.1\%. A change in the frequency with which a learner predicts a class can affect its precision simultaneously. The appendix retains all classes and all groups, so these individual changes are not used to reinterpret the overall macro-F1 effect.

The validation results also explain why the final test was necessary. During development the TCN GAF-versus-observed-only difference was $-1.33$ points and the GRU difference was $+0.90$ points. Both paired intervals spanned zero. On the final cohort the TCN decline was larger and the GRU point estimate changed sign. A claim of successful augmentation based solely on the positive validation GRU estimate would therefore have failed even as a point prediction of the final comparison.

Under uniform sampling, each observed training window has expected exposure $73{,}600/1{,}437$ in the observed-only group and $73{,}600/1{,}862$ in an augmented group. These are roughly 51.2 and 39.5 presentations, respectively. Thus a fixed training budget involves a concrete exchange: additional examples replace a fraction of the presentations of observed windows. The observed-only result should be interpreted at that budget. Holding original-example exposure constant would examine a different perturbation and would require a new experimental protocol.
